\documentclass[../../main.tex]{subfiles}
\begin{document}

\begin{flushright}
\footnotesize\textit{【自動文字起こし・要確認】}
\end{flushright}

\noindent
{\bf [解]}
\paragraph{(1)}
\begin{align}
1\text{回シャッフル}&\ \cdots\ \{5,1,6,2,7,3,8,4\}\\
2\text{回シャッフル}&\ \cdots\ \{7,5,3,1,8,6,4,2\}\\
3\text{回シャッフル}&\ \cdots\ \{8,7,6,5,4,3,2,1\}
\end{align}
から，$\{8,7,6,5,4,3,2,1\}$．

\paragraph{(2)}

$\{1,2,\ldots,2N\}$を$t$回シャッフルした後における数$k$（$1\le k\le2N$）の位置を$f_t(k)$と表す．以下の漸化式を得る．
\begin{align}
1\le f_t(k)\le N\text{の時，}&\ f_{t+1}(k)=2f_t(k)\\
N+1\le f_t(k)\le2N\text{の時，}&\ f_{t+1}(k)=2f_t(k)-(2N+1)
\end{align}
$\cdots①$

（シャッフルでは，現在の並びの前半（位置$1,\ldots,N$）の$i$番目の数は新しい並びの位置$2i$に，後半（位置$N+1,\ldots,2N$）の$i$番目の数は新しい並びの位置$2i-1$に移ることによる．）

題意から，$f_0(k)=k$，$f_1(k)=f(k)$であるから，①より，
\begin{align}
1\le k\le N\text{の時，}&\ f(k)-2k=(2k-2k)=0\\
N+1\le k\le2N\text{の時，}&\ f(k)-2k=(2k-2k)-(2N+1)=-(2N+1)
\end{align}
となり，いずれも$f(k)-2k$は$2N+1$で割り切れる．$\blacksquare$

\paragraph{(3)}

以下合同式の法を$2N+1=2^n+1$とする．①から，$1\le f_t(k)\le N$，$N+1\le f_t(k)\le2N$いずれの場合も
\begin{align}
f_{t+1}(k)\equiv2f_t(k)\pmod{2N+1}
\end{align}
（わざわざ(2)を使えば，$f(k)\to f_{t+1}(k)$，$k\to f_t(k)$として同じ式を得る．）だから，くり返し用いて$f_0(k)=k$より，
\begin{align}
f_t(k)\equiv2^t\cdot k\pmod{2N+1}
\end{align}

だから，$t=2n$として
\begin{align}
f_{2n}(k)\equiv2^{2n}k=\left(2^n\right)^2k\equiv\left((2^n+1)-1\right)^2k\equiv(-1)^2k=k\pmod{2N+1}
\end{align}

ここで，$1\le k\le2N$，$1\le f_{2n}(k)\le2N$だから，法が$2N+1$であること（区間の幅$2N$が法より小さいこと）より，合同から等号が従い，
\begin{align}
f_{2n}(k)=k
\end{align}

したがって，題意は示された．$\blacksquare$

\end{document}
